% Part XVIII: The Jungerman-Ringel Theorem
\section{Part XVIII: The Jungerman--Ringel Theorem Is $W(3,3)$ in Action}

The 1980 paper of Jungerman and Ringel~\cite{JR1980} determines the minimum number $\delta(S_p)$ of triangles in a minimal triangulation of every closed orientable surface~$S_p$ of genus~$p$.  We show that \emph{every structural element} of that theorem is parameterised by $W(3,3)$.

\subsection{The Master Formula in $W(3,3)$ Notation}

Jungerman--Ringel Theorem 1.1 states:
\[
  \delta(S_p) \;=\; 2\!\left\lceil\frac{\Phi_6 + \sqrt{1 + \mu k\,p}}{2}\right\rceil + \mu(p - 1),
  \qquad p \neq \lambda,
\]
with the unique exception $\delta(S_\lambda) = f = 24$.  Here $\Phi_6 = 7$, $\mu = 4$, $k = 12$, $\lambda = 2$, $f = 24$ are the standard $W(3,3)$ parameters.  The equivalent Theorem~1.2 asserts that every admissible pair $(n,t)$ satisfying
\[
  4 \le n,\quad 0 \le t \le n - q!,\quad (n-q)(n-q-1) \equiv 2t \pmod{k}
\]
admits a triangular embedding---except the unique pair $(q^2, q) = (9,3)$.

\subsection{The Twelve Cases = $k$ Residue Classes}

The proof splits into twelve cases by $n \bmod k$.  The four \emph{allowed} residue classes, where $K_n$ itself triangulates, are
\[
  \{0,\; q,\; \mu,\; \Phi_6\} = \{0, 3, 4, 7\} \pmod{k}.
\]
These are the four fundamental $W(3,3)$ parameters modulo the valence.  The eight \emph{forbidden} classes require edge removal and handle subtraction.

\medskip\noindent
The \textbf{current graph index} organises the twelve classes into a gauge-theoretic decomposition:

\begin{center}
\begin{tabular}{ccl}
\hline
Index & Residues mod $k$ & Algebraic role \\ \hline
$1$ & $\{0, q, \mu, \Phi_6\}$ & Full symmetry: direct $K_n$ embedding \\
$2$ & $\{2, 6, 8, 10\}$ & Chiral: $\mathbb{Z}_2$ quotient, principle C10 \\
$3$ & $\{1, 5, 9\}$ & Color: $\mathbb{Z}_3$ quotient, principle C7 \\
var & $\{11\}$ & Exceptional: $K_n - K_{q+\lambda}$ sector \\
\hline
\end{tabular}
\end{center}

\noindent
The splitting $4 + 4 + 3 + 1 = k = 12$ matches the gauge boson count of $\mathrm{SU}(3)\times\mathrm{SU}(2)\times\mathrm{U}(1)$:
\begin{itemize}
\item Index~$1$ ($4$ classes): electroweak sector $\mathrm{SU}(2)_L \times \mathrm{U}(1)_Y$;
\item Index~$3$ ($3$ classes $\to$ $8$ gluons): color sector $\mathrm{SU}(3)_c$;
\item Index~$2$ ($4$ classes): chiral mixing (parity violation);
\item Exceptional ($1$ class): $\mathrm{SU}(5)$ GUT sector ($K_n - K_{q+\lambda}$).
\end{itemize}

\subsection{Handle Subtraction = Chain Complex Boundary}

Lemma~3.1 of~\cite{JR1980} defines \emph{handle subtraction}: replacing a hexagonal pattern in the Heffter scheme removes $q! = 6$ edges while preserving Rule~R*.  In the dual:
\[
  \Delta t = q!, \qquad \Delta(\text{dual vertices}) = -\mu, \qquad \Delta(\text{genus}) = -1.
\]
This is the boundary operator $\partial$ of the chain complex
\[
\mathrm{Points}(v) \to \mathrm{Pairs}(45) \to \mathrm{Spreads}(27),
\]
quantised in units of $q!$.
The \emph{arithmetic comb} provides $m$ independent handles simultaneously, creating a \textbf{harmonic oscillator} with unit level spacing:
\begin{align*}
  K_k &\colon \lambda = 2\text{ oscillator levels}, \\
  K_f &\colon \mu = 4\text{ levels}, \\
  K_v &\colon q! = 6\text{ levels}.
\end{align*}
The level count is $n/q!$ for $n = k, f$ and $\lfloor(v-q!)/q!\rfloor + 1 = q!$ for $n = v$.

\subsection{Lock~13: The Unique Jungerman--Ringel Obstruction}

The pair $(q^2, q) = (9, 3)$ is the \emph{unique} obstruction in the orientable theorem.  We prove this selects $q = 3$:

\begin{theorem}[Lock~13]
Among all prime powers $q$, the congruence $(q^2 - 3)(q^2 - 4) \equiv 2q \pmod{12}$ holds only when $q \equiv 3 \pmod{6}$, i.e.\ $q \in \{3, 9, 15, \ldots\}$.  Moreover, $q^2$ lands in a forbidden residue class ($q^2 \equiv 9 \bmod 12$) only for $q \equiv 3 \pmod{6}$.  Since $q = 3$ is the smallest such prime power and the only one small enough for the combinatorial obstruction to hold (Huneke's proof requires exhaustive enumeration at $n = q^2 = 9$), the Jungerman--Ringel exception selects $q = 3$.
\end{theorem}

The resolution uses $(n,t) = (\Phi_4, q^2) = (10, 9)$, giving $\delta(S_\lambda) = f = 24$.  The gap $f - \chi = 24 - 22 = \lambda$ equals the mass ratio $k/q!$.

\subsection{Ternary Induction = Generation Functor}

Theorem~4.10.1 of~\cite{JR1980} is a \emph{ternary induction}: from triangular embeddings of $K_{3t-1} - h_1$, $K_{3t} - h_2$, $K_{3t+1} - h_3$ (three consecutive integers), one constructs an embedding of $K_{3t+q} - (h_1 + h_2 + h_3 + q)$.  The step size is $q = 3$, and the three inputs combine into one output---the \textbf{ternary generation functor} $\mathcal{F}_q$.  Applied repeatedly, this builds all Case~9 constructions from $10$ base cases, precisely paralleling three fermion generations combining under the $q = 3$ structure.

\subsection{The Discriminant Boolean Cube}

The discriminant $1 + \mu k\,p$ is a perfect square exactly when $K_n$ triangulates directly (the Heawood bound is tight).  At $W(3,3)$ parameters:
\[
  \sqrt{1 + 48 \cdot \mathrm{genus}(K_n)} = 2n - \Phi_6 \quad\text{for } n \in \{\mu, \Phi_6, k, g, f, q^3, v\}.
\]
In particular $\sqrt{1 + 48 \times 111} = 73 = \Phi_{12}$, so the twelfth cyclotomic value is the discriminant root at $n = v$.

The perfect-square residues modulo $\mu k = 48$ form a group $(\mathbb{Z}/2\mathbb{Z})^q$ of order $2^q = 8$, the \textbf{Boolean cube} of dimension~$q$.  Its four positive generators $\{1, \Phi_6, k{+}q{+}\lambda, f{-}1\} = \{1, 7, 17, 23\}$ form a Klein four-group under multiplication mod~$48$.

Mapping the cube's three $\mathbb{Z}_2$ factors to isospin, hypercharge, and color gives $2^q = 8$ fermion types per generation, whence $q \times 2^q = 3 \times 8 = f = 24$ Weyl fermions total---exactly the $f = 24$ triangles on the exceptional double torus $S_\lambda$.

\subsection{The Dual Polyhedra Tower}

The tower of minimal triangulations at small genus gives a ladder of self-dual polyhedra:
\begin{center}
\begin{tabular}{ccccl}
\hline
$h$ & $v$ & $e$ & $f$ & Name \\ \hline
$0$ & $\mu = 4$ & $6$ & $4$ & Tetrahedron (sphere) \\
$1$ & $\Phi_6 = 7$ & $21$ & $14$ & Cs\'asz\'ar (torus) \\
$\lambda = 2$ & $\Phi_4 = 10$ & $36$ & $f = 24$ & JR resolution (double torus) \\
$q! = 6$ & $k = 12$ & $66$ & $44$ & Heffter $K_{12}$ (genus~$q!$) \\
\hline
\end{tabular}
\end{center}

\noindent
The vertex--Jordan pairing links each level to a division algebra:
\[
  \mu + \dim J_3(\mathbb{R}) = \Phi_4, \qquad
  \Phi_6 + \dim J_3(\mathbb{H}) = \chi, \qquad
  \Phi_4 + \dim J_3(\mathbb{O}) = v - q.
\]
The differences between consecutive sums are $k$ and $g$---themselves $W(3,3)$ parameters.

\subsection{Kirchhoff's Law = Gauge Invariance}

The current graph axioms encode gauge theory:
\begin{itemize}
\item Principle C4 (Kirchhoff's current law at non-vortex vertices): $\nabla\cdot J = 0$, the conservation law of gauge invariance.
\item Vortex vertices violate C4 by an amount that \emph{generates} the cyclic group---they are \textbf{charged particles}.
\item Principle C7 (index~$3$ subgroup): the quotient $G/H \cong \mathbb{Z}_3$ is colour triality.
\item Principle C10 (index~$2$ subgroup): the quotient $G/H \cong \mathbb{Z}_2$ is chirality (left/right).
\end{itemize}

The $10$ construction principles C1--C10 can be viewed as the axioms of a \textbf{topological gauge theory} on the surface, with $q = 3$ controlling the maximum valence (C1), the colour quotient (C7), and the relationship between vortex charge and group generation (C5).

\subsection{Summary: Lock Count}

With Lock~13 (Jungerman--Ringel obstruction) added to the previous twelve, the total selection pressure stands at:
\[
  \boxed{13\text{ independent locks, all selecting } q = 3.}
\]

\begin{thebibliography}{9}
\bibitem{JR1980} M.~Jungerman and G.~Ringel, ``Minimal triangulations on orientable surfaces,'' \textit{Acta Math.}\ \textbf{145} (1980), 121--154.
\end{thebibliography}
